\section{Directional Derivatives}\label{sec:directional_derivative}

Partial derivatives give us an understanding of how a surface changes when we move in the $x$ and $y$ directions. We made the comparison to standing in a rolling meadow and heading due east: the amount of rise/fall in doing so is comparable to $f_x$. Likewise, the rise/fall in  moving due north is comparable to $f_y$. The steeper the slope, the greater in magnitude $f_y$.

But what if we didn't move due north or east? What if we needed to move northeast and wanted to measure the amount of rise/fall? Partial derivatives alone cannot measure this. This section investigates \textbf{directional derivatives}, which do measure this rate of change. 

We begin with a definition.

\begin{definition}[Directional Derivatives]\label{def:direct_deriv}%
Let $z=f(x,y)$ be continuous on an open set $S$ and let $\vec u =\bracket{u_1,u_2}$ be a unit vector. For all points $(x,y)$, the \textbf{directional derivative of $f$ at $(x,y)$ in the direction of $\vec u$} is
\index{derivative!directional}\index{directional derivative}
\[D_{\vec u\,}f(x,y) = \lim_{h\to 0} \frac{f(x+hu_1,y+hu_2) - f(x,y)}h.\]
\end{definition}

The partial derivatives $f_x$ and $f_y$ are defined with similar limits, but only $x$ or $y$ varies with $h$, not both. Here both $x$ and $y$ vary with a weighted $h$, determined by a particular unit vector $\vec u$. This may look a bit intimidating but in reality it is not too difficult to deal with; it often just requires extra algebra. However, the following theorem reduces this algebraic load.

% ``at'' was ``on an open set $S$ containing''
\begin{theorem}[Directional Derivatives]\label{thm:direct_deriv1}%
Let $z=f(x,y)$ be differentiable at $(x_0,y_0)$, and let $\vec u =\bracket{u_1,u_2}$ be a unit vector. The directional derivative of $f$ at $(x_0,y_0)$ in the direction of $\vec u$ is
\index{derivative!directional}\index{directional derivative}
\[D_{\vec u\,}f(x_0,y_0)=f_x(x_0,y_0)u_1 + f_y(x_0,y_0)u_2.\]
\end{theorem}

%\begin{proof}
%This is really a quick application of \autoref{def:multi_differentiability}.  Because $f$ is differentiable at $(x_0,y_0)$,
%\begin{align*}
% D_{\vec u\,}f(x_0,y_0)
% &=\lim_{h\to0}
% \frac{f_x(x_0,y_0)hu_1+f_y(x_0,y_0)hu_2+E_x(x_0,y_0)hu_1+E_y(x_0,y_0)hu_2}h \\
% &=f_x(x_0,y_0)u_1 + f_y(x_0,y_0)u_2+\lim_{h\to0}E_x(x_0,y_0)u_1+E_y(x_0,y_0)u_2 \\
% &=f_x(x_0,y_0)u_1 + f_y(x_0,y_0)u_2.\qedhere
%\end{align*}
%\end{proof}

\begin{proof}
We will use \autoref{thm:multi_chain}, with $x=g(t)=u_1t+x_0$ and $y=h(t)=u_2t+y_0$.  This means that
\begin{align*}
D_{\vec u\,}f(x_0,y_0)
&=\frac{\dd f}{\dd t}(x_0,y_0)\\
&=f_x(x_0,y_0)\frac{\dd x}{\dd t}(0)+f_y(x_0,y_0)\frac{\dd y}{\dd t}(0)\\
&=f_x(x_0,y_0)u_1 + f_y(x_0,y_0)u_2.\qedhere
\end{align*}
\end{proof}

\youtubeVideo{uCY0XYXgYQo}{Finding the Directional Derivative --- Ex 1}

\begin{example}[Computing directional derivatives]\label{ex_direct1}%
Let $z= 14-x^2-y^2$ and let $P=(1,2)$. Find the directional derivative of $f$, at $P$, in the following directions:
\begin{enumerate}
	\item toward the point $Q=(3,4)$,
	\item	in the direction of $\bracket{2,-1}$, and
	\item toward the origin.
\end{enumerate}
\solution
\mtable{Understanding the directional derivative in \autoref{ex_direct1}.}{fig:direct1}{%
\myincludeasythree{
3Droll=-0.45935076561172045,
3Dortho=0.00494238780811429,
3Dc2c=0.8379672765731812 0.26980355381965637 0.4743594527244568,
3Dcoo=-27.334510803222656 46.63825988769531 23.82315444946289,
3Droo=130.00000000067573}{\marginparwidth}{At P, unit vector u in the xy‑plane gives directional derivative D_u f = nabla f(P)* u; nabla f arrow drawn on the sheet.}{figures/figdirect1_3D}}
The surface is plotted in \autoref{fig:direct1}, where the point $P=(1,2)$ is indicated in the $x,y$-plane as well as the point $(1,2,9)$ which lies on the surface of $f$. We find that $f_x(x,y) = -2x$ and $f_x(1,2) = -2$; $f_y(x,y) = -2y$ and $f_y(1,2) = -4$. 
\begin{enumerate}
	\item Let $\vec u_1$ be the unit vector that points from the point $(1,2)$ to the point $Q=(3,4)$, as shown in the figure. The vector $\vv{PQ} =\bracket{2,2}$; the unit vector in this direction is $\vec u_1=\bracket{1/\sqrt{2}, 1/\sqrt{2}}$. Thus the directional derivative of $f$ at $(1,2)$ in the direction of $\vec u_1$ is
	\[D_{\vec u_1}f(1,2) = -2(1/\sqrt{2}) +(-4)(1/\sqrt{2}) = -6/\sqrt{2}%\approx -4.24
	.\]
	Thus the instantaneous rate of change in moving from the point $(1,2,9)$ on the surface in the direction of $\vec u_1$ (which points toward the point $Q$) is $-3\sqrt2$. Moving in this direction moves one steeply downward.
	
	\item		We seek the directional derivative in the direction of $\bracket{2,-1}$. The unit vector in this direction is $\vec u_2 =\bracket{2/\sqrt{5},-1/\sqrt{5}}$. Thus the directional derivative of $f$ at $(1,2)$ in the direction of $\vec u_2$ is
	\[D_{\vec u_2}f(1,2) = -2(2/\sqrt{5})+(-4)(-1/\sqrt{5}) = 0.\]
	Starting on the surface of $f$ at $(1,2)$ and moving in the direction of $\bracket{2,-1}$ (or $\vec u_2$) results in no instantaneous change in $z$-value. This is analogous to standing on the side of a hill and choosing a direction to walk that does not change the elevation. One neither walks up nor down, rather just ``along the side'' of the hill.
	
	Finding these directions of ``no elevation change'' is important.
	
	\item		At $P=(1,2)$, the direction towards the origin is given by the vector $\bracket{-1,-2}$; the unit vector in this direction is $\vec u_3=\bracket{-1/\sqrt{5},-2/\sqrt{5}}$. The directional derivative of $f$ at $P$ in the direction of the origin is
	\[D_{\vec u_3}f(1,2) = -2(-1/\sqrt{5}) + (-4)(-2/\sqrt{5}) = 10/\sqrt{5}% \approx 4.47
	.\]
	Moving towards the origin means ``walking uphill'' quite steeply, with an initial slope of $2\sqrt5$.
\end{enumerate}
\end{example}

As we study directional derivatives, it will help to make an important connection between the unit vector $\vec u =\bracket{u_1,u_2}$ that describes the direction and the partial derivatives $f_x$ and $f_y$. We start with a definition and follow this with a Key Idea.

\begin{definition}[Gradient]\label{def:gradient}%
Let $z=f(x,y)$ be differentiable on an open set $S$ that contains the point $(x_0,y_0)$.
\index{gradient} 
\begin{enumerate}
	\item The \textbf{gradient of $f$} is $\grad f(x,y) =\bracket{f_x(x,y),f_y(x,y)}$.
	\item The \textbf{gradient of $f$ at} $(x_0,y_0)$ is $\grad f(x_0,y_0) =\bracket{f_x(x_0,y_0),f_y(x_0,y_0)}$.
\end{enumerate}
\end{definition}

\mnote{\textbf{Note:} The symbol ``$\nabla$'' is named ``nabla,'' derived from the Greek name of a Jewish harp. Oddly enough, in mathematics the expression $\grad f$ is pronounced ``del $f$.''}

To simplify notation, we often express the gradient as $\grad f =\bracket{f_x, f_y}$.
%It is often useful to think of the gradient $\grad$ as an operator:\index{del operator}
%\[
%\grad=\bracket{\frac\partial{\partial x},\frac\partial{\partial y}}.
%\]
%The operator $\grad$ only has any meaning when it operates on a function; it doesn't mean anything by itself.  But this notation does help us to apply it correctly to find the gradient.
The gradient allows us to compute directional derivatives in terms of a dot product.

\begin{keyidea}[The Gradient and Directional Derivatives]\label{idea:gradient_direct}%
%Let $z=f(x,y)$ be differentiable on an open set $S$ and let $\vec u$ be a unit vector. 
The directional derivative of $z=f(x,y)$ in the direction of the unit vector $\vec u$ is
\index{gradient}\index{directional derivative}\index{derivative!directional}
\[D_{\vec u\,}f = \grad f\cdot \vec u.\]
\end{keyidea}

The properties of the dot product previously studied allow us to investigate the properties of the directional derivative. Given that the directional derivative gives the instantaneous rate of change of $z$ when moving in the direction of $\vec u$, three questions naturally arise:
\begin{enumerate}
	\item In what direction(s) is the change in $z$ the greatest (i.e., the ``steepest uphill'')?
	\item In what direction(s) is the change in $z$ the least (i.e.,  the ``steepest downhill'')?
	\item In what direction(s) is there no change in $z$?
\end{enumerate}

Using the key property of the dot product, we have
\begin{equation}
\grad f\cdot \vec u = \norm{\grad f}\,\vnorm u \cos \theta = \norm{\grad f}\cos \theta, \label{eq:gradient}
\end{equation}
where $\theta$ is the angle between the gradient and $\vec u$. (Since $\vec u$ is a unit vector, $\vnorm{u} = 1$.) This equation allows us to answer the three questions stated previously.
\begin{enumerate}
	\item \autoeqref{eq:gradient} is maximized when $\cos \theta =1$, i.e., when the gradient and $\vec u$ have the same direction. We conclude the gradient points in the direction of greatest $z$ change.
	\item	\autoeqref{eq:gradient} is minimized when $\cos \theta = -1$, i.e., when the gradient and $\vec u$ have opposite directions. We conclude the gradient points in the opposite direction of the least $z$ change.
	\item \autoeqref{eq:gradient} is 0 when $\cos \theta = 0$, i.e., when the gradient and $\vec u$ are orthogonal to each other. We conclude the gradient is orthogonal to directions of no $z$ change. 
\end{enumerate}

This result is rather amazing. Once again imagine standing in a rolling meadow and face the  direction that leads you steepest uphill. Then the direction that leads steepest downhill is directly behind you, and side-stepping either left or right (i.e., moving perpendicularly to the direction you face) does not change your elevation at all.

%Recall that a level curve is defined by a path in the $x$-$y$ plane along which the $z$-values of a function do not change; the directional derivative in the direction of a level curve is 0. This is analogous to walking along a path in the rolling meadow along which the elevation does not change. The gradient at a point is orthogonal to the direction where the $z$ does not change; i.e., the gradient is orthogonal to level curves.

Recall that a level curve is defined as a curve in the $x$-$y$ plane along which the $z$-values of a function do not change. Let a surface $z=f(x,y)$ be given, and let's represent one such level curve as a vector-valued function, $\vrt =\bracket{x(t), y(t)}$. As the output of $f$ does not change along this curve, $f\bigl(x(t),y(t)\bigr) = c$ for all $t$, for some constant $c$.

Since $f$ is constant for all $t$, $\frac{\dd f}{\dd t} = 0$. By the Multivariable Chain Rule, we also know
\begin{align*}
\frac{\dd f}{\dd t} &= f_x(x,y)x\primeskip'(t) + f_y(x,y)y\primeskip'(t)\\
	&=\bracket{f_x(x,y),f_y(x,y)}\cdot\bracket{x\primeskip'(t),y\primeskip'(t)}\\
						&= \grad f \cdot \vrp(t)\\
						&=0.
\end{align*}
This last equality states $\grad f \cdot \vrp(t) = 0$: the gradient is orthogonal to the derivative of $\vec r$. Our conclusion: at any point on a surface, the gradient at that point is orthogonal to the level curve that passes through that point.

We restate these ideas in a theorem, then use them in an example.

\begin{theorem}[The Gradient and Directional Derivatives]\label{thm:gradient}%
Let $z=f(x,y)$ be differentiable on an open set $S$ with gradient $\grad f$, let $P=(x_0,y_0)$ be a point in $S$ and let $\vec u$ be a unit vector.
\index{gradient}\index{directional derivative}\index{derivative!directional}\index{level curves}\index{gradient!and level curves}
\begin{enumerate}
	\item The maximum value of $D_{\vec u\,}f(x_0,y_0)$ is $\norm{\grad f(x_0,y_0)}$; the direction of maximal $z$ increase is $\grad f(x_0,y_0)$.
	%obtained when the angle between $\grad f$ and $\vec u$ is 0, i.e.,  the direction of maximal increase is $\grad f$.
	\item   The minimum value of $D_{\vec u\,}f(x_0,y_0)$ is $-\norm{\grad f(x_0,y_0)}$; the direction of minimal $z$ increase is $-\grad f(x_0,y_0)$.
	%The minimum value of $D_{\vec u\,}f$ is $-\norm{\grad f}$, obtained when the angle between $\grad f$ and $\vec u$ is $\pi$, i.e., the direction of minimal increase is $-\grad f$.
	\item At $P$, $\grad f(x_0,y_0)$ is orthogonal to the level curve passing through $(x_0,y_0)$.
	%$D_{\vec u\,}f = 0$ when $\grad f$ and $\vec u$ are orthogonal.
\end{enumerate}
\end{theorem}

\begin{example}[Finding directions of maximal and minimal increase]\label{ex_direct2}%
Let $f(x,y) = \sin x\cos y$ and let $P=(\pi/3,\pi/3)$. Find the directions of maximal/minimal increase, and find a direction where the instantaneous rate of $z$ change is 0.
\solution
We begin by finding the gradient. We see that the partial derivatives are $f_x = \cos x\cos y$ and $f_y = -\sin x\sin y$, so that
\[\grad f =\bracket{\cos x\cos y,-\sin x\sin y}\quad \text{and, at $P$,} \quad \grad f\left(\frac{\pi}3,\frac{\pi}3\right) =\bracket{\frac14,-\frac34}.\]
Thus the direction of maximal increase is $\bracket{1/4, -3/4}$. In this direction, the instantaneous rate of $z$ change is $\norm{\bracket{1/4,-3/4}} = \sqrt{10}/4% \approx 0.79
$. 

\mtable{Graphing the surface and important directions in \autoref{ex_direct2}.}{fig:direct2}{%
\myincludeasythree{
3Droll=0,
3Dortho=0.0051408205181360245,
3Dc2c=0.6666666865348816 0.6666666865348816 0.3333333432674408,
3Dcoo=-1.1194640398025513 -3.250971794128418 8.740827560424805,
3Droo=150}{.8\marginparwidth}{Level set f=c (blue ellipse) sits on z = f; black nabla f(P) is normal, pointing in steepest‑ascent direction.}{figures/figdirect2_3D}
\\(a)\\
\myincludeasythree{
3Droll=0,
3Dortho=0.008940433152019978,
3Dc2c=0.6666666865348816 0.6666666865348816 0.3333333134651184,
3Dcoo=-6.498621940612793 -16.487354278564453 45.97212600708008,
3Droo=149.9999987284343}{.8\marginparwidth}{Edge view: nabla f is vertical; red tangent line satisfies nabla f*T = 0, confirming orthogonality to the gradient.}{figures/figdirect2b_3D}
\\(b)}

\autoref{fig:direct2} shows the surface plotted from two different perspectives. In each, the gradient is drawn at $P$ with a dashed line (because of the nature of this surface, the gradient points ``into'' the surface). Let $\vec u =\bracket{u_1, u_2}$ be the unit vector in the direction of $\grad f$ at $P$. Each graph of the figure also contains the vector $\bracket{u_1, u_2, \norm{\grad f\,}}$. This vector has a ``run'' of 1 (because in the $x$-$y$ plane it moves 1 unit) and a ``rise'' of $\norm{\grad f}$, hence we can think of it as a vector with slope of $\norm{\grad f}$ in the direction of $\grad f$, helping us visualize how ``steep'' the surface is in its steepest direction. 
%To help understand this number, consider moving from $P$ one unit in the direction of $\bracket{1/4,-3/4}$. The unit vector in this direction is $\bracket{1/\sqrt{10},-3/\sqrt{10}}\approx\bracket{0.316,-0.949}$; moving from $P$ one unit in the direction of the gradient places one at the point $R=(1.363,0.098)$. 

%Compare the $z$-values at $P$ and $R$:
%\[f\left(\frac{\pi}3,\frac{\pi}3\right) = \frac{\sqrt{3}}4 \approx 0.433.\]

The direction of minimal increase is $\bracket{-1/4,3/4}$; in this direction the instantaneous rate of $z$ change is $-\sqrt{10}/4% \approx -0.79
$.

Any direction orthogonal to $\grad f$ is a direction of no $z$ change. We have two choices: the direction of $\bracket{3,1}$ and the direction of $\bracket{-3,-1}$. The unit vector in the direction of $\bracket{3,1}$ is shown in each graph of the figure as well. The level curve at $z=\sqrt{3}/4$ is drawn: recall that along this curve the $z$-values do not change. Since $\bracket{3,1}$ is a direction of no $z$-change, this vector is tangent to the level curve at $P$.
\end{example}

\begin{example}[Understanding when $\grad f = \vec 0$]\label{ex_direct9}%
Let $f(x,y) = -x^2+2x-y^2+2y+1$. Find the directional derivative of $f$ in any direction at $P=(1,1)$.
\solution
We find $\grad f =\bracket{-2x+2, -2y+2}$. At $P$, we have $\grad f(1,1) =\bracket{0,0}$. 
According to \autoref{thm:gradient}, this is the direction of maximal increase. However, $\bracket{0,0}$ is directionless; it has no displacement. And regardless of the unit vector $\vec u$ chosen, $D_{\vec u\,}f = 0$.

\mtable{At the top of a paraboloid, all directional derivatives are 0.}{fig:direct9}{%
\myincludeasythree{
3Droll=0.641848137510025,
3Dortho=0.00471429293975234,
3Dc2c=0.8800624012947083 0.28546664118766785 0.3794717788696289,
3Dcoo=-34.139766693115234 45.736881256103516 38.29633712768555,
3Droo=200.00000572048617}{\marginparwidth}{Paraboloid peak of z = x^2 + y^2: nabla f = 0, hence D_u f = 0 for every horizontal direction u.}{figures/figdirect9_3D}}

\autoref{fig:direct9} helps us understand what this means. We can see that $P$ lies at the top of a paraboloid. In all directions, the instantaneous rate of change is 0. 

So what is the direction of maximal increase? It is fine to give an answer of $\vec 0 =\bracket{0,0}$, as this indicates that all directional derivatives are 0.
\end{example}

The fact that the gradient of a surface always points in the direction of steepest increase/decrease is very useful, as illustrated in the following example.

\begin{example}[The flow of water downhill]\label{ex_direct3}%
Consider the surface given by $f(x,y)= 20-x^2-2y^2$. Water is poured on the surface at $(1,1/4)$. What path does it take as it flows downhill?
\solution
Let $\vrt =\bracket{x(t), y(t)}$ be the vector-valued function describing the path of the water in the $x$-$y$ plane; we seek $x(t)$ and $y(t)$. We know that water will always flow downhill in the steepest direction; therefore, at any point on its path, it will be moving in the direction of $-\grad f$. (We ignore the physical effects of momentum on the water.) Thus $\vrp(t)$ will be parallel to $\grad f$, and there is some constant $c$ such that $c\grad f = \vrp(t) =\bracket{x\primeskip'(t), y\primeskip'(t)}$. 

We find $\grad f =\bracket{-2x, -4y}$ and write $x\primeskip'(t)$ as $\frac{\dd x}{\dd t}$ and $y\primeskip'(t)$ as $\frac{\dd y}{\dd t}$. Then 
\begin{align*}
c\grad f &=\bracket{x\primeskip'(t), y\primeskip'(t)}\\
\bracket{-2cx, -4cy}& =\bracket{\frac{\dd x}{\dd t}, \frac{\dd y}{\dd t}}.
\end{align*}
This implies
\begin{gather*}
-2cx = \frac{\dd x}{\dd t} \quad \text{and} \quad  -4cy =\frac{\dd y}{\dd t}, \text{ i.e.,}\\
c = -\frac{1}{2x}\frac{\dd x}{\dd t} \quad \text{and} \quad  c =-\frac{1}{4y}\frac{\dd y}{\dd t}.
\end{gather*}
As $c$ equals both expressions, we have
\[\frac{1}{2x}\frac{\dd x}{\dd t} =\frac{1}{4y}\frac{\dd y}{\dd t}.\]
To find an explicit relationship between $x$ and $y$, we can integrate both sides with respect to $t$. Recall from our study of differentials that $\ds \frac{\dd x}{\dd t}\dd t = \dd x$. Thus:
\begin{align*}
\int \frac{1}{2x}\frac{\dd x}{\dd t}\dd t &=\int \frac{1}{4y}\frac{\dd y}{\dd t}\dd t \\
\int \frac{1}{2x}\dd x &=\int\frac{1}{4y}\dd y \\
\frac 12\ln\abs x &= \frac14\ln\abs y +C_1\\
2\ln\abs x &= \ln\abs y +C_1\\
\ln \abs{x^2} &= \ln\abs y+C_1% &= y,
\intertext{Now raise both sides as a power of \intertextmath{e}:}
%\end{align*}
%where we skip some algebra in the last step. 
%\begin{align*}
x^2 &= e^{\ln \abs y+C_1}\\
x^2 &= e^{\ln \abs y}e^{C_1}\qquad \text{(Note that $e^{C_1}$ is just a constant.)}\\
x^2 &= yC_2\\
\frac1{C_2}x^2 &=y   \qquad \text{(Note that $1/C_2$ is just a constant.)}\\
Cx^2 &= y.
\end{align*}
As the water started at the point $(1,1/4)$, we can solve for $C$:
%
\mtable{A graph of the surface described in \autoref{ex_direct3} along with the path in the $x$-$y$ plane with the level curves.}{fig:direct3}{%
\myincludeasythree{
3Droll=0.03187335123623557,
3Dortho=0.004501350689679384,
3Dc2c=0.37703937292099 0.8674991130828857 0.3244788944721222,
3Dcoo=31.472492218017578 -33.2381477355957 46.63370132446289,
3Droo=150.00001038122664}{\marginparwidth}{Curve r(t) ascends a parabolic shell z = g(x,y); the tangent arrow shows dr/dt whose df/dt = nabla f·r′.}{figures/figdirect3_3D}
\\(a)\\
\begin{tikzpicture}[alt={Concentric red level curves f=c; blue spiral path cuts them, picturing how f grows along r(t).}]
\begin{axis}[width=1.16\marginparwidth,tick label style={font=\scriptsize},
axis y line=middle,axis x line=middle,name=myplot,
ymin=-3.5,ymax=3.5,xmin=-4.2,xmax=4.2]
\addplot[thick,draw={\colortwo},smooth,domain=0:360]({2*cos(x)},{sqrt(2)*sin(x)});
\addplot[thick,draw={\colortwo},smooth,domain=0:360]({2*sqrt(2)*cos(x)},{2*sin(x)});
\addplot[thick,draw={\colortwo},smooth,domain=0:360]({2*sqrt(3)*cos(x)},{sqrt(6)*sin(x)});
\addplot[thick,draw={\colortwo},smooth,domain=0:360]({4*cos(x)},{2*sqrt(2)*sin(x)});
\addplot[thick,draw={\colorone},smooth,domain=1:3.1]({x},{x^2/4});
\end{axis}
\node [right] at (myplot.right of origin) {\scriptsize $x$};
\node [above] at (myplot.above origin) {\scriptsize $y$};
\end{tikzpicture}
\\(b)}%
%
\[C(1)^2 = \frac14 \quad \Rightarrow \quad C = \frac14.\]
Thus the water follows the curve $y=x^2/4$ in the $x$-$y$ plane. The surface and the path of the water is graphed in \autoref{fig:direct3}(a). In part (b) of the figure, the level curves of the surface are plotted in the $x$-$y$ plane, along with the curve $y=x^2/4$. Notice how the path intersects the level curves at right angles. As the path follows the gradient downhill, this reinforces the fact that the gradient is orthogonal to level curves.
\end{example}

\subsection{Functions of Three Variables}

The concepts of directional derivatives and the gradient are easily extended to three (and more) variables. We combine the concepts behind Definitions \ref{def:direct_deriv} and \ref{def:gradient} and \autoref{thm:direct_deriv1} into one set of definitions.

\begin{definition}[Directional Derivatives and Gradient with Three Variables]\label{def:direct_deriv3}%
Let $w=F(x,y,z)$ be differentiable on an open ball $B$ and let $\vec u $ be a unit vector in $\mathbb{R}^3$.\index{gradient}\index{directional derivative}\index{derivative!directional}
\begin{enumerate}
	\item	The \textbf{gradient} of $F$ is $\grad F =\bracket{F_x,F_y,F_z}$.
	\item The \textbf{gradient of $F$ at $(x_0,y_0,z_0)$ is
	\[\grad F(x_0,y_0,z_0) =\bracket{f_x(x_0,y_0,z_0),f_y(x_0,y_0,z_0),f_z(x_0,y_0,z_0)}.\]}
	\item The \textbf{directional derivative of $F$ in the direction of $\vec u$} is
	\[D_{\vec u\,}F=\grad F\cdot \vec u.\]
\end{enumerate}
\end{definition}

The same properties of the gradient given in \autoref{thm:gradient}, when $f$ is a function of two variables, hold for $F$, a function of three variables.

{\tcbset{grow to right by=10em}
\begin{theorem}[The Gradient and Directional Derivatives with Three Variables]\label{thm:gradient3}%
Let $w=F(x,y,z)$ be differentiable on an open set $S$ with gradient $\grad F$, let $P=(x_0,y_0,z_0)$ be a point in $S$, and let $\vec u$ be a unit vector.\index{gradient}\index{directional derivative}\index{derivative!directional}
\begin{enumerate}
	\item The maximum value of $D_{\vec u\,}F(x_0,y_0,z_0)$ is $\norm{\grad F(x_0,y_0,z_0)}$; %, obtained when the angle between $\grad F$ and $\vec u$ is 0, i.e.,
	 the direction of maximal increase is $\grad F(x_0,y_0,z_0)$.
	\item The minimum value of $D_{\vec u\,}F(x_0,y_0,z_0)$ is $-\norm{\grad F(x_0,y_0,z_0)}$; %, obtained when the angle between $\grad F$ and $\vec u$ is $\pi$, i.e.,
	 the direction of minimal increase is $-\grad F(x_0,y_0,z_0)$.
	\item At $P$, $D_{\vec u\,}F(x_0,y_0,z_0) = 0$ when $\grad F(x_0,y_0,z_0)$ and $\vec u$ are orthogonal.
%	\item At $P$, $\grad F(x_0,y_0,z_0)$ is orthogonal to the level surface passing through
\end{enumerate}
\end{theorem}}

We interpret the third statement of the theorem as ``the gradient is orthogonal to level surfaces,'' the three-variable analogue to level curves.
\index{gradient!and level surfaces}\index{level surface}

\begin{example}[Finding directional derivatives with functions of three variables]\label{ex_direct5}%
If a point source $S$ is radiating energy, the intensity $I$ at a given point $P$ in space is inversely proportional to the square of the distance between $S$ and $P$. That is, when $S=(0,0,0)$,  $\ds I(x,y,z) = \frac{k}{x^2+y^2+z^2}$ for some constant $k$.

Let $k=1$, let $\vec u =\bracket{2/3, 2/3, 1/3}$ be a unit vector, and let $P = (2,5,3).$ Measure distances in inches. Find the directional derivative of $I$ at $P$ in the direction of $\vec u$, and find the direction of greatest intensity increase at $P$.
\solution
We need the gradient $\grad I$, meaning we need $I_x$, $I_y$ and $I_z$. Each partial derivative requires a simple application of the Quotient Rule, giving
\begin{align*}
\grad I &=\bracket{\frac{-2x}{(x^2+y^2+z^2)^2},\frac{-2y}{(x^2+y^2+z^2)^2},\frac{-2z}{(x^2+y^2+z^2)^2}}\\
\grad I(2,5,3) &=\bracket{\frac{-4}{1444},\frac{-10}{1444},\frac{-6}{1444}}%\approx\bracket{-0.003,-0.007,-0.004}
\\
D_{\vec u\,}I &= \grad I(2,5,3)\cdot \vec u\\
&= -\frac{17}{2166}% \approx -0.0078
.
\end{align*}
The directional derivative tells us that moving in the direction of $\vec u$ from $P$ results in a slight decrease in intensity% of about $-0.008$ units per inch
. (The intensity is decreasing as $\vec u$ moves one farther from the origin than $P$.)

The gradient gives the direction of greatest intensity increase. Notice that 
\begin{align*}
\grad I(2,5,3) &=\bracket{\frac{-4}{1444},\frac{-10}{1444},\frac{-6}{1444}}\\
			&= \frac{2}{1444}\bracket{-2,-5,-3}.
\end{align*}
That is, the gradient at $(2,5,3)$ is pointing in the direction of $\bracket{-2,-5,-3}$, that is, towards the origin. That should make intuitive sense: the greatest increase in intensity is found by moving towards to source of the energy.
\end{example}

The directional derivative allows us to find the instantaneous rate of $z$ change in any direction at a point. We can use these instantaneous rates of change to define lines and planes that are \emph{tangent} to a surface at a point, which is the topic of the next section.

\printexercises{exercises/12-05-exercises}
